\chapter{Introduction}
\label{ch:introduction}

\section{Motivation}

Modern derivative engines must do more than price options: they must calibrate models repeatedly and produce stable sensitivities for risk management. In the Heston stochastic-volatility model \cite{heston}, semi-analytical and Fourier-based methods provide a strong numerical foundation because the characteristic function is available. In particular, the Fourier-COS method gives accurate and efficient reference prices when the truncation and expansion settings are controlled \cite{fang_oosterlee_2009}.

The practical difficulty is that these routines are rarely called only once. Calibration requires repeated evaluations of the pricing map, large benchmark datasets require many reference valuations, and Greek analysis requires stable local derivatives. This repeated-evaluation setting makes neural approximators attractive, but speed alone is not sufficient in a financial workflow. A learned component is useful only if its errors are measurable, its failure regions are visible, and its outputs remain compatible with the downstream task.

This thesis studies neural methods from that perspective. The question is not whether neural networks can be fast, but whether they can be made reliable inside a benchmarked derivative workflow. The objective is therefore to identify where neural components can accelerate Heston pricing, calibration, and Greek computation while remaining anchored to explicit numerical references and clearly stated limitations.

\section{Problem Setting}

The thesis focuses on European vanilla options under the Heston model. This setting is narrow enough to admit reliable reference methods, but rich enough to expose the main difficulties of neural derivative analytics: nonlinear parameter effects, repeated inverse problems, low-price regions, short maturities, and non-smooth payoff behavior.

The main methodological risk is that average pricing accuracy can hide downstream failures. A surrogate with small mean pricing error may still distort calibration if its residuals are structured across moneyness or maturity. A differentiable pricing network may produce plausible prices while giving unstable Delta or Gamma. These issues are especially relevant near the short-maturity at-the-money region, where the payoff kink creates a structural challenge for second-order sensitivities.

For this reason, the thesis treats benchmarking as part of the method rather than as an afterthought. Numerical prices, implied-volatility labels, calibrated parameters, PINN prices, and Greeks are all evaluated under explicit reference conventions. This makes it possible to separate acceleration from reliability, and to distinguish useful neural components from model variants that only improve a single metric in isolation.

\section{Methodological Approach}

The workflow combines a numerical reference lane with a neural approximation lane. The reference lane produces controlled Heston-COS prices, implied-volatility labels, and semi-analytical Greek benchmarks. The neural lane uses COS/IV labels for the ANN surrogate, embeds the surrogate in a CaNN calibration loop, and evaluates PINN pricing and Greek components against the reference objects. The base PINN is trained through the PDE formulation; reference Greeks enter only through benchmark diagnostics and the Sobolev-enhanced variant.

\begin{figure}[H]
    \centering
    \resizebox{0.98\textwidth}{!}{%
    \begin{tikzpicture}[
        font=\small,
        box/.style={draw, rounded corners=2pt, align=center, minimum height=0.9cm, text width=2.65cm, inner sep=4pt},
        refbox/.style={box, fill=blue!7},
        nnbox/.style={box, fill=green!8},
        outbox/.style={box, fill=gray!10},
        arrow/.style={-{Latex[length=2.5mm]}, thick},
        evalarrow/.style={-{Latex[length=2.5mm]}, thick, dashed},
        node distance=1.05cm and 1.05cm
    ]
        \node[refbox] (setup) {Heston setup\\synthetic contracts};
        \node[refbox, right=of setup] (cos) {COS pricing\\IV inversion};

        \node[nnbox, below=1.0cm of cos] (ann) {ANN IV\\surrogate};
        \node[nnbox, right=of ann] (cann) {CaNN\\calibration};
        \node[nnbox, right=of cann] (pinn) {PINN pricing\\surface};
        \node[nnbox, right=of pinn, yshift=0.55cm] (sobolev) {Sobolev\\training};
        \node[nnbox, right=of pinn, yshift=-0.55cm] (acv) {ACV\\control variate};
        \node[outbox, right=of sobolev, yshift=-0.55cm] (greeks) {Greek\\diagnostics};
        \node[refbox, above=1.0cm of greeks] (refs) {Reference\\Greeks};

        \node[align=right, text width=2.4cm, left=0.55cm of setup] (reflabel) {Numerical\\reference};
        \node[align=right, text width=2.4cm, left=0.55cm of ann] (nnlabel) {Neural\\workflow};

        \draw[arrow] (setup) -- (cos);
        \draw[arrow] (cos) -- (refs);
        \draw[arrow] (cos) -- (ann);
        \draw[arrow] (ann) -- (cann);
        \draw[arrow] (cann) -- (pinn);
        \draw[arrow] (pinn.east) -- (sobolev.west);
        \draw[arrow] (pinn.east) -- (acv.west);
        \draw[arrow] (sobolev.east) -- (greeks.north west);
        \draw[arrow] (acv.east) -- (greeks.south west);
        \draw[evalarrow] (refs.south) -- (greeks.north);
        \draw[evalarrow] (refs.west) -- (sobolev.north);
    \end{tikzpicture}}
    \caption{High-level structure of the benchmarked Heston workflow. Solid arrows show the main computational flow; dashed arrows show reference Greeks used to compute diagnostic errors and to define Sobolev Greek targets.}
    \label{fig:intro_pipeline}
\end{figure}

This organization is important for the interpretation of the results. The ANN is judged not only as a fast function approximator, but also as the forward engine used by the calibration routine. The PINN is judged not only by price RMSE, but also by spatial error structure and derivative quality. The Greek analysis is treated as a separate problem because automatic differentiation through a learned price surface does not by itself guarantee reliable sensitivities.

\section{Main Contributions}

The first contribution is a reproducible benchmark pipeline for Heston derivative analytics. The pipeline connects data generation, COS reference pricing, implied-volatility inversion, ANN training, CaNN calibration, PINN pricing, and Greek diagnostics under common conventions.

The second contribution is the improved neural calibration result. Building on the Liu-style calibration framework \cite{liu}, the final \texttt{tanh} ANN/CaNN configuration improves the controlled Liu et al. benchmark in the reported calibration setting, especially in parameter recovery.

The third contribution is a PINN pricing and sensitivity benchmark under Heston. The PINN is evaluated against COS and semi-analytical references not only at the price level, but also through first- and second-order Greeks.

The fourth contribution is the analysis of Greek improvement methods in the hard short-maturity at-the-money region. Sobolev-style derivative supervision \cite{czarnecki2017sobolev} and the analytic-control-variate diagnostic improve Greek quality, while the remaining errors show that the payoff kink remains a structural limitation.

\section{Scope and Thesis Structure}

The scope is intentionally controlled. The experiments use European vanilla options under the Heston model, synthetic data with known reference values, and benchmark protocols where numerical settings can be inspected. Market microstructure effects, bid-ask spreads, dividends, exotic payoffs, and portfolio-scale deployment are outside the main experimental scope. This restriction is deliberate: it allows the thesis to test the numerical and neural components before making broader claims about market deployment.

The remainder of the thesis is organized as follows. Chapter~\ref{ch:background} introduces the mathematical and financial background, including the Heston model, characteristic-function pricing, the COS method, semi-analytical Greeks, and implied-volatility inversion. Chapter~\ref{ch:ann_surrogate} presents the supervised ANN implied-volatility surrogate. Chapter~\ref{ch:cann} describes the CaNN calibration workflow. Chapter~\ref{ch:pinn_greeks} develops the PINN pricing and Greek computation framework. Chapter~\ref{ch:performance_benchmarks} reports the benchmark results for ANN pricing, CaNN calibration, PINN pricing, Sobolev Greeks, and the ACV diagnostic. Chapter~\ref{ch:conclusions} summarizes the conclusions, limitations, and future work.
